% Generated from validated Article Draft Result. Do not hand-edit; revise the structured draft instead.
\section{総括――2024 H1に何が変わったのか}
\label{pkg:conclusion}
\noindent\textbf{2024年上期を振り返ると、生成AIの競争単位は「最も強い一つのモデル」から、interaction、openness、deployment、architecture、action、media、control、servingを組み合わせるsystem stackへ広がった。各章の出来事を再列挙せず、その構造変化を一本の線として結び直す。}\autocite{src-054d1f81750b,src-0fe343211040,src-2d3a533e91db,src-4ff5ab7d45bc,src-77e35e5f22be,src-b5fbfc0d43e4,src-c89d38f0e830,src-f9616f30e8f5}

2024-H1の最大の変化は、model qualityが重要でなくなったことではない。むしろ高性能化が進んだ結果、model単体のbenchmarkだけでは実用上の差を説明できなくなったことである。GPT-4oのinteraction、Llama 3に象徴されるopen-weight ecosystem、Appleのon-device/server配置、Jambaのhybrid architecture、Command Rのretrieval/tool接続、Soraの生成メディア、Anthropicのinterpretability、NVIDIA NIMのserving――これらは別々の流行ではなく、能力の成立条件がstack全体へ拡張した同じ構造変化の断面である。\autocite{src-054d1f81750b,src-0fe343211040,src-2d3a533e91db,src-4ff5ab7d45bc,src-77e35e5f22be,src-b5fbfc0d43e4,src-c89d38f0e830,src-f9616f30e8f5}


各layerは相互に依存する。multimodal interactionを広げればlatencyとservingが効き、open modelを配備すればlicense・runtime・data provenanceが問題になり、edgeへ移せばmemoryとmodel sizeが制約になる。toolやcomputer environmentへ接続すればbehavior policyとsecurity boundaryが重くなり、生成メディアの表現力が上がればavailabilityや公開条件の違いが実用性を左右する。2024-H1は、個々の技術が一本のroadmapへ収束したのではなく、複数layerを組み合わせて初めて「使えるAI」を定義する時代への入口だった。\autocite{src-054d1f81750b,src-0fe343211040,src-2d3a533e91db,src-4ff5ab7d45bc,src-77e35e5f22be,src-c89d38f0e830,src-f9616f30e8f5}


2024年6月末までに一次情報から確定できるのは、長文脈・multimodality・open model・small/on-device・hybrid/sparse architecture・RAG/tool use・video/image generation・control/interpretability・serving packagingが、それぞれ独立した競争軸として明確になったということである。一方で、どのarchitectureが長期的に主流になるか、agentが実運用でどこまで安定するか、openの定義がどこまでtraining transparencyを含むか、生成videoが制作workflowへどの速度で定着するかまではH1の資料だけでは確定できない。\autocite{src-054d1f81750b,src-0fe343211040,src-2d3a533e91db,src-4ff5ab7d45bc,src-77e35e5f22be,src-b5fbfc0d43e4,src-c89d38f0e830,src-f9616f30e8f5}


したがってH2へ持ち越された問いは、「次にどのmodelが首位になるか」だけではない。reasoningとactionをどう接続するか、open modelをどう効率よく運用するか、localとcloudをどう分担するか、multimodal systemの権限と安全をどう管理するか、生成メディアをpreviewからproductionへどう移すか、そしてmodel更新の速度にevaluation・runtime・controlが追随できるかである。これは後から得たH2の答えをH1へ投影するのではなく、H1末にすでに見えていた未解決条件として残すべき問いである。\autocite{src-054d1f81750b,src-0fe343211040,src-2d3a533e91db,src-4ff5ab7d45bc,src-77e35e5f22be,src-b5fbfc0d43e4,src-c89d38f0e830,src-f9616f30e8f5}


\begin{claimboundary}
本総括は、2024-H1の一次資料と検証済みEvidenceから導くretrospective synthesisである。vendor/project/authorの評価値を独立再現済みの事実へ格上げせず、2026年から見た後知恵で2024年6月末時点のavailabilityや成熟度を書き換えない。H2との接続は、当時残っていた問いとしてのみ記述する。\autocite{src-054d1f81750b,src-0fe343211040,src-2d3a533e91db,src-2f7e28810d05,src-4ff5ab7d45bc,src-5019b9256d83,src-77e35e5f22be,src-b5fbfc0d43e4,src-c89d38f0e830,src-edd253654067,src-f9616f30e8f5}
\end{claimboundary}
